\section{Additional Home Sales Results}
\label{sec:appendix_sales_rental}

This appendix presents additional results for home sale prices, complementing the event study analysis in Section~\ref{sec:empirical_results_home_sales}.

\subsection{Event Study Without Hedonic Controls}

The main text presents event study results including property-level hedonic controls to isolate the pure price effect holding quality constant. Figure~\ref{fig:event_study_sales_no_hedonics} presents results without hedonic controls to capture the total effect of aldermanic strictness on housing costs, including any composition effects operating through changes in the housing stock. Point estimates are similar, suggesting that composition effects account for only a modest portion of the total price response.

\begin{figure}[H]
    \centering
    \begin{minipage}{0.48\textwidth}
        \centering
        \includegraphics[width=\textwidth]{../tasks/run_event_study_sales_disaggregate/output/event_study_disaggregate_yearly_cohort_2015_continuous_triangular_1000ft_no_hedonics_full.pdf}
        \caption*{Panel A: Continuous Treatment}
    \end{minipage}
    \hfill
    \begin{minipage}{0.48\textwidth}
        \centering
        \includegraphics[width=\textwidth]{../tasks/run_event_study_sales_disaggregate/output/event_study_disaggregate_yearly_cohort_2015_continuous_split_triangular_1000ft_no_hedonics_full.pdf}
        \caption*{Panel B: Continuous Split}
    \end{minipage}
    \caption{Event Study: Home Sale Prices --- Without Hedonic Controls}
    \figurenotes{Same specification as Figure~\ref{fig:event_study_sales} but excluding property-level hedonic controls. 2015 redistricting cohort. Border-pair-by-side and border-pair-by-year fixed effects. Triangular kernel weights with 1,000 ft bandwidth. Year $-1$ omitted as reference. 95\% confidence intervals based on standard errors clustered at the census block level.}
    \label{fig:event_study_sales_no_hedonics}
\end{figure}

\subsection{Alternative Event Timing: Announcement vs.\ Implementation}

The main specification defines treatment timing based on when redistricting takes effect (2015 implementation). If housing markets are forward-looking, prices may begin responding when redistricting maps are first announced (2012). Figure~\ref{fig:event_study_sales_announcement} presents the event study using 2012 announcement timing. The pattern is qualitatively similar, with point estimates slightly larger under announcement timing.

\begin{figure}[H]
    \centering
    \includegraphics[width=0.6\textwidth]{../tasks/run_event_study_sales_disaggregate/output/event_study_disaggregate_yearly_stacked_announcement_continuous_triangular_1000ft_full.pdf}
    \caption{Event Study: Home Sale Prices --- Announcement Timing}
    \figurenotes{Same specification as Figure~\ref{fig:event_study_sales} but using 2012 announcement timing rather than 2015 implementation timing. Triangular kernel weights with 1,000 ft bandwidth. Year $-1$ omitted as reference. 95\% confidence intervals based on standard errors clustered at the census block level.}
    \label{fig:event_study_sales_announcement}
\end{figure}